\documentclass[a4paper]{article}

%% Language and font encodings
\usepackage[english]{babel}
\usepackage[utf8x]{inputenc}
% \usepackage[T1]{fontenc}

%% Sets page size and margins
\usepackage[a4paper,top=3cm,bottom=2cm,left=3cm,right=3cm,marginparwidth=1.75cm]{geometry}

%% Useful packages
\usepackage{amsmath}
\usepackage{graphicx}
\usepackage[colorlinks=true, allcolors=blue]{hyperref}
\usepackage{apacite} % APA style citations
\AtBeginDocument{\urlstyle{APACsame}}  % Links in APA citytions same formatting
\usepackage{tabulary}
\usepackage{natbib} % natbib citations: \citep{} and \citet{} for in-text
\usepackage{booktabs}
\usepackage{float}
\usepackage{adjustbox}
\usepackage[skip=0.5ex]{subcaption}
% If you use \tableofcontents, this adjusts the name
%\addto\captionsenglish{\renewcommand*\contentsname{Table of Contents}}
\newcommand{\ubar}[1]{\text{\b{$#1$}}}


\title{\textbf{Spatial Road Network Misallocation and Optimal Expansion in Europe and Central Asia}}
\author{Sebastian Krantz and Ana Cusolito}

\begin{document}
\maketitle

This brief examines whether road infrastructure was allocated efficiently across European and Central-Asian (ECA) countries, using a quantitative spatial general equilibrium model developed by \citet{fajgelbaum2020optimal} and taking inspiration from \citet{graff2024spatial} for its calibration with the Open Street Map (OSM) routing engine (OSRM) \citep{luxen-vetter-2011}. As in \citet{fajgelbaum2020optimal}, I consider a 0.5$^\circ$ grid of estimates with GDP per capita and population as a starting point for constructing national network representations---using recent global estimates by \citet{LGEAW}. I then determine the population centroid of each cell using 1$km^2$ estimates for the year 2020 from WorldPop \citep{tatem2017worldpop}. I retain all cells covering the continental ECA region, excluding island states like the UK, the nordic peninsula (Norway, Sweden, Finland), and limiting Russia to the north by St. Petersburg and to the east by Novosibirsk. Thus, I obtain a grid of 7694 connected 0.5$^\circ$ cells covering 43 ECA countries. \newline 

Each cell touches up to 8 neighbouring cells. I then retrieve fastest car routes between the population centroids of all neighbouring cells using OSRM. The median route has a length of 93km and spans a geodesic distance of 58.2$km$. Since OSRM 'snaps' the centroids to the nearest road, I follow \citet{graff2024spatial} and add the geodesic distance from the centroids to the route starting/end points to the route length, assuming a travel speed of 10$km/h$ to not inflate travel times too much. The median travel speed on the routes themselves is 47.6$km/h$, and adding the start/end segments at 10$km/h$ reduces the median speed to 42.2$km/h$. \newline 

 Overall infrastructure efficiency/quality is thus a function of two factors: (1) route efficiency (RE), which is the geodesic distance divided by the total length of the route from centroid to centroid; and (2) travel speed along the total route. I combine these two factors by calculating time efficiency (TE) as the geodesic centroids distance divided by the total travel time. It effectively measures the travel speed in the direction of the destination, accounting for both infrastructure 
 
 \begin{figure}[H] % \vspace{-3mm}
 \centering
 \caption{\label{fig:ECAgnet} Time Efficiency in 0.5$^\circ$ Grid Network Spanning Continental ECA}
 \includegraphics[width=\textwidth, trim= {5mm 0 5mm 0}, clip]{"../figures/ECA_grid_network_time_efficiency.pdf"} %trim={<left> <lower> <right> <upper>}
 \vspace{-5mm}
 \end{figure}
 
\noindent quality and quantity connecting two cells. The median TE estimate is 23.6$km/h$. Figure \ref{fig:ECAgnet} shows the result. As expected, infrastructure is much more efficient in Western Europe where some segments have TE of up to 82$km/h$. It is also notable that TE along the Russion border are near zero---due to the Ukraine war and other geopolitical tensions (e.g. with Poland) resulting in border restrictions taken into account in OSM. This will be less of a practical issue as in the following I will only assess national infrastructure optimality using the model of \citet{fajgelbaum2020optimal}. \newline 
 
Concretely, for each country, I retain all national centroids and centroids within 100$km$ geodesic distance from national centroids, alongside all connecting segments. I then let the 15 largest cells with a population above 200,000 produce different products which are traded across the network. Other cells are divided into 5 quantiles according to population, where each quantile produces its own product. Thus, in total 20 products are traded and each cell only produces one of them. This amounts to an \citet{armington1969theory} model of products differentiated by origin, and a recent meta-study by \citet{BAJZIK2020103383} finds a median Armington Elasticity (of substitution) of $\sigma = 3.8$, which I accordingly set in \citet{fajgelbaum2020optimal}'s model. The goods are produced using a linear technology $Y_i = Z_i L_i$, where $L_i$ is cell $i$'s population and $Z_i$ cell GDP per Capita in 2017 PPP terms---both obtained from  \citet{LGEAW}. Moreover, individual utility over traded and non-traded goods (housing) is a Cobb-Douglas function $U = c^\alpha h^{1-\alpha}$, where following \citet{fajgelbaum2020optimal} I set $\alpha = 0.4$, and following \citet{graff2024spatial} I set $h = 1$, i.e., housing is proportional to population and each person is given one unit of housing. I also follow \citet{fajgelbaum2020optimal} by setting the congestion parameter $\beta = 0.13$, and the returns to infrastructure $\gamma = 0.1$ in the baseline (decreasing returns to infrastructure) case, and I set iceberg trade costs to $\delta^\tau_{ij} = 0.116 \times \log(\text{geodesic distance in km})$ roughly following \citet{graff2024spatial}. These parameters enter an iceberg trade costs equation of the form $1+\tau^n_{ij}$ capturing the cost of trading good $n$ between cells $i$ and $j$, with $\tau^n_{ij} = \delta^\tau_{ij} (Q^n_{ij})^\beta / I_{ij}^\gamma.$ I also consider the increasing returns to infrastructure case by inverting the ratio of $\beta$ and $\gamma$ following \citet{fajgelbaum2020optimal} to obtain $\gamma' = \beta^2/\gamma = 0.169$. \newline 

The latter may be more realistic since \citet{shepherd2007trade}, studying trade and road infrastructure in ECA, find that a 1\% improvement in average road quality is associated with a 0.2\% to 0.6\% increase in trade, and \citet{cocsar2016domestic}, studying large-scale public investment in roads undertaken in Turkey during the 2000s, find that the distance elasticity of trade improved by 0.10–0.27 when roads were upgraded. Both studies thus suggest a value of $\gamma$ well above 0.1. \newline 

Finally, following \citet{fajgelbaum2020optimal}, and \citet{collier2016cost}, and \citet{krantz2024optimal}, I estimate the cost of building infrastructure on a link as 
\begin{equation}
\delta^I_{ij} = \delta^I_0 + 0.12 \times \log(\text{Ruggedness}_{ij}) + 0.085 \times \log(\text{Population per $km^2$}_{ij}),
\end{equation}
where terrain Ruggedness is taken from \citet{nunn2012ruggedness} and averaged in each cell, and population density is computed using WorldPop 2020 estimates \citep{tatem2017worldpop}. To obtain bilateral measures, I average the estimates of the respective cells $i$ and $j$. The network building constraint is then $K = \sum_i \sum_j \delta^I_{ij} I_{ij}$, and, following \citet{fajgelbaum2020optimal}, I determine $\delta^I_0$ such that the budget $K = $1,000,000. In the following I then run two simulations for each country following \citet{graff2024spatial}: (1) an optimal 10\% expansion of the existing infrastructure stock ($I_{ij} = $ time efficiency) where $K = $1,100,000 and $\ubar{I}_{ij} = I_{ij}$ (existing infrastructure is fixed) and $\bar{I}_{ij} = 70km/h$, which is in the 99.9th percentile; and (2) an optimal reallocation exercise where $K = $1,000,000 and $\ubar{I}_{ij} = \min(I_{ij}, 10km/h)$, i.e., the welfare maximizing social planner can reallocate infrastructure. \newline 

In optimizing infrastructure, the planner only maximizes total national welfare, and can also only spend/reallocate infrastructure on national links. The 100$km$ buffer zone takes surrounding cells into account as trading partners. Infrastructure may thus be built to foster international trade, but only within national boundaries, and only if it maximizes welfare for domestic consumers. % \newline 

\subsection*{Optimal Reallocation}
 
Figure \ref{fig:Realloc} shows the welfare gains from the optimal reallocation of national networks of 39 ECA countries, excluding Russia where no solution could be obtained due to the very large network. Under increasing returns ($\gamma = 0.169 > \beta$), the gains are higher due to overall lower trade costs $\tau_{ij}$ resulting in larger volumes of trade in this case. But the gains under DR ($\gamma = 0.1 < \beta$) are 

 \begin{figure}[H] \vspace{-10mm}
 \centering
 \caption{\label{fig:Realloc} Welfare Gains from Optimal Reallocation of National Road Networks}
 \vspace{-3mm}
 % \includegraphics[width=\textwidth]{"../figures/ECA_grid_network_realloc_barchart_both_decomp.pdf"} %trim={<left> <lower> <right> <upper>}
 \includegraphics[width=\textwidth]{"../figures/ECA_grid_network_realloc_barchart_both_decomp_nograd.pdf"} %trim={<left> <lower> <right> <upper>}
 \\ \vspace{3mm}
 \begin{adjustbox}{center}
 \begin{tabular}{cc}
 Estonia (0.68\%) & Kyrgyztan (0.64\%) \\
 \includegraphics[width=0.6\textwidth, trim = {1cm 5cm 1cm 5cm}, clip]{"../figures/grid_network/realloc_fixed_irs/EST_opt_realloc_irs.pdf"} &  \includegraphics[width=0.6\textwidth, trim = {1cm 5cm 1cm 5cm}, clip]{"../figures/grid_network/realloc_fixed_irs/KGZ_opt_realloc_irs.pdf"} \\ 
 Tajikistan (0.55\%) & Bulgaria (0.53\%) \\
 \includegraphics[width=0.6\textwidth, trim = {1cm 5cm 1cm 4cm}, clip]{"../figures/grid_network/realloc_fixed_irs/TJK_opt_realloc_irs.pdf"} &  \includegraphics[width=0.6\textwidth, trim = {1cm 5cm 1cm 5cm}, clip]{"../figures/grid_network/realloc_fixed_irs/BGR_opt_realloc_irs.pdf"} \\ 
 \end{tabular}
 \end{adjustbox}
 \end{figure}
 
\noindent very similar between the different countries. In both cases, there are four countries: Bulgaria, Tajikistan, Kyrgyzstan, and Estonia, with the highest welfare gains. \newline 

The bottom half of Figure \ref{fig:Realloc} therefore shows the graphical results for these countries. The thickness of links amounts to the final reallocation, whereas their colour indicates whether infrastruc-ture was added (red) or removed (blue) during the reallocation process. The size of nodes corresponds to their population, and their colour indicates welfare gains from reallocation, with red or orange nodes experiencing gains, green nodes maintaining their welfare level, and blue nodes loosing out on welfare. National nodes surrounded by a white circle or buffer nodes surrounded by a black circle produce their own unique product and are thus particularly important trading partners. \newline 

The overall reallocation pattern exhibited by these four countries, as well as most other countries, is that infrastructure is taken out of peripheral and sparsely populated regions to better connect urban centers. The results thus lend support to the hypothesis that, at least in terms of road connectivity, too much infrastructure spending targeted unimportant roads, whereas larger welfare gains could have been achieved from improving inter-city connectivity. Yet, in all cases, the welfare gains are below 1\%, indicating that road networks in ECA are generally more efficient than, e.g., in Africa, where some countries like Somalia can reap gains of $\geq$6\% from reallocation \citep{graff2024spatial}. \newline 

 The lighter shaded bars in Figure \ref{fig:Realloc} further visualize the change in welfare attributable to a volume effect measured as the total increase in consumption applied proportionally to all locations. The remaining welfare change is due to an allocation effect. In the median country, 86\% of the welfare gain is explained by the volume effect, with mild heterogeneity across countries ranging from 61\%-100\%. Thus, a more optimal allocation of road infrastructure would have generated welfare gains predominantly through increased trade volumes between major cities.  \newline 
 
We can also seek to better understand the misallocation by correlating it with various country characteristics such as total land area ($r = 0.31$) GDP per Capita ($r = -0.23$), the Logistics Performance Index ($r = -0.2$), the component of it measuring 'Quality of Trade and Transport-Related Infrastructure' ($r = -0.18$), and the 2020 Doing Business Score ($r = -0.19$). Thus, misallocation is mildly correlated with development indicators and total land area, suggesting that both the capacity to plan and the planning scope/economic dispersement affect outcome optimality. \newline 

Appendix Figure \ref{fig:ReallocExtra} shows the reallocation graphs for further countries under increasing returns ($\gamma > \beta$) in order of their welfare gains. As indicated, the general pattern is that infrastructure is reallocated from the periphery to segments connecting urban centers, as this is where most people live and thus where most welfare gains can be reaped. A concave individual utility function alone is thus not sufficient to incentivize a welfare maximizing social planner to spend significantly on infrastructure in scarcely populated rural regions. Other planning objectives may correct for this. % \newline 

\subsection*{Optimal 10\% Expansion}

Analogous, Figure \ref{fig:10pExp} visualizes the welfare gains from an optimal 10\% network expansion. The gains are lower than from reallocation, indicating that even an optimal a 10\% investment is not able to equal the benefits from a more optimal past resource allocation. Again, Estonia would reap the largest gains of around 0.35\%, and the ranking of countries is overall similar to the reallocation case. It is more notably, however, that some western European countries such as Greece and Denmark could also reap large gains from a 10\% optimal network expansion. This may be correlated with their their more fragmented geographies and the resulting difficulties of building efficient integrated road networks on them. Similar to the reallocation case, optimal expansions mostly target major links connecting urban centers, complemented by small peripheral investments. \newline 

Appendix Figure \ref{fig:10PExpExtra} again shows the optimal expansions for further countries in order of their welfare gains. In most countries, the investments are also allocated along central links, with some supporting investments to 'feeder lines'. \newline

The analysis overall thus suggests that European and particularly Central Asian countries could reap moderate welfare gains from an optimal expansion of their road networks targeting mainly

 \begin{figure}[H] \vspace{-10mm}
 \centering
 \caption{\label{fig:10pExp} Welfare Gains from Optimal 10\% Expansion of National Road Networks}
  \vspace{-3mm}
 % \includegraphics[width=\textwidth]{"../figures/ECA_grid_network_10perc_barchart_both.pdf"} %trim={<left> <lower> <right> <upper>}
 \includegraphics[width=\textwidth]{"../figures/ECA_grid_network_10perc_barchart_both_decomp_nograd.pdf"}
 \\ \vspace{3mm}
 \begin{adjustbox}{center}
 \begin{tabular}{cc}
 Estonia (0.35\%) & Greece (0.25\%) \\
 \includegraphics[width=0.6\textwidth, trim = {1cm 5cm 1cm 7cm}, clip]{"../figures/grid_network/10perc_fixed_irs/EST_opt_10perc_irs.pdf"} & \includegraphics[width=0.6\textwidth, trim = {1cm 5cm 1cm 4cm}, clip]{"../figures/grid_network/10perc_fixed_irs/GRC_opt_10perc_irs.pdf"}  \\ 
 Tajikistan (0.21\%) & Kyrgyztan (0.19\%) \\
 \includegraphics[width=0.6\textwidth, trim = {1cm 5cm 1cm 4cm}, clip]{"../figures/grid_network/10perc_fixed_irs/TJK_opt_10perc_irs.pdf"} & \includegraphics[width=0.6\textwidth, trim = {1cm 5cm 1cm 5cm}, clip]{"../figures/grid_network/10perc_fixed_irs/KGZ_opt_10perc_irs.pdf"}  \\ 
 \end{tabular}
 \end{adjustbox}
 \end{figure}

\noindent central links connecting major urban centers. This is not a consequence of the IR ($\gamma > \beta$) assumption leading to more tree-like networks as described by \citet{fajgelbaum2020optimal}; the results under DR ($\gamma < \beta$) are very similar and thus not shown. 

\subsection*{Expansion Matching Reallocation Gains}

Rather than computing an ad-hoc 10\% optimal network expansion and evaluating the welfare gains of to each country, we can also approximately solve for the network expansion needed to match the welfare gains from an optimal reallocation of the existing network. Figure \ref{fig:MRealloc} shows the results, which are surprisingly similar across the IRS and DRS (black bars) scenarios, ranging from a 17.3\% to a 40.8\% network expansion needed to reach the gains from optimal reallocation.

 \begin{figure}[H] % \vspace{-1mm}
 \centering
 \caption{\label{fig:MRealloc} Increase in Infrastructure Stock to Match Optimal Reallocation Gains}
 \vspace{-3mm}
 % \includegraphics[width=\textwidth]{"../figures/ECA_grid_network_realloc_barchart_both_decomp.pdf"} %trim={<left> <lower> <right> <upper>}
 \includegraphics[width=\textwidth]{"../figures/ECA_grid_network_Mrealloc_barchart_both_nograd.pdf"} %trim={<left> <lower> <right> <upper>}
 \end{figure}

 
The correlation with the optimal reallocation gains is 0.42 under IRS and 0.31 under DRS, indicating that countries with a less optimal allocation of infrastructure also need to spend more to 'fix it'. In the median country, the infrastructure investment needed is 81.4 times larger than the welfare gains from reallocation (both in percentage terms), indicating that it is costly to fix welfare losses from past misallocation with new investments. 



\subsection*{Optimal Investments in ECA's Transport Network}

Apart from national network inefficiencies, I also consider the transport network of the entire macro-region and simulate optimal investments into it. To do this, I first select the largest cities with a population above 250,000 that are at least 50$km$ apart from each other, using data from the Simplemaps World Cities Database \citep{simplemaps} built from authoritative sources. This yields 281 key cities. I then use OSRM to obtain fastest car routes between them, and intersect and process their geometries following \citet{krantz2024optimal} to obtain an undirected graph representation of the continental road network comprising 952 nodes and 1722 edges. I assign to each node city populations within 30$km$ of that node, and calibrate its productivity as the inverse-distance weighted GDP per capita in 2017 PPP terms across all cell-centroids within 50$km$ of the node using the 0.25$^\circ$ data by \citet{LGEAW}. I also reparameterize each edge with the road distance and travel time from OSRM. Furthermore, I apply a network finding algorithm following \citet{krantz2024optimal} to the network to find high-potential new/missing links in the network and verify that they can be constructed. This yields 274 potential links, highlighted in green in Figure \ref{fig:ECAnet}, which would decrease road distance by at least 50\% along the nodes it connects. 


\begin{figure}[H]
\centering
\caption{\label{fig:ECAnet} ECA's Transport Network Graph}
\includegraphics[width=\textwidth, trim = {0.5cm 0.5cm 0.5cm 0.5cm}, clip]{"../figures/trans_ECA_network_actual_discretized_gravity_new_roads_real_edges.pdf"}
\scriptsize 
\emph{Notes:} Figure shows ECA transport network connecting 281 key cities (blue dots), including proposed links in green.
\end{figure}

The average route efficiency across existing links is 0.8, and I assume that the same holds for proposed links. Thus, their length is assumed $1/0.8 = 1.25$ times longer than the geodesic distance spanned by them. Figure \ref{fig:ECAnet} also visualizes a normalized sum of gravity along the existing links, which is the sum, across all routes using a particular link of the product of origin and destination city populations divided by the road distance. Thus, it is a measure of the expected traffic along the link, assuming open borders. In this regard, I also note that the Ukraine-Russia border is presently closed and OSRM does not route across it. Consequently, roads crossing it are not part of the existing network, but the network finding algorithm found several links reconnecting Russia to Ukraine. Since much physical infrastructure was destroyed in the war, it may be sensible to assume that these links need to be rebuilt, although at a lower cost than new roads. \newline 

Speaking of which, I also estimate the cost of constructing/upgrading each link following \citet{krantz2024optimal} using an equation of the form
\begin{equation} \label{eq:COST}
\ln\left(\frac{\text{cost}}{km}\right) = \ln(X) -0.11 (\text{dist} > 50km) + 0.12 \ln(\text{rugg}) + 0.085 \ln\left(\frac{\text{pop}}{km^2}\right),
\end{equation}
where $X = $100,000 for new links and lower for different types of upgrades following \citet{krantz2024optimal} and information in the ROCKS database \citep{rocks2018}. Appendix Figure \ref{fig:ECAcalibdet} visualizes various network calibration details discussed above. Ostensibly, the network is Western Europe is much better and faster than in Eastern Europe, Russia, and Central Asia. 


\subsubsection*{Optimal Investments in Partial Equilibrium}

To analyze the cost-benefit of building/upgrading each link, I simulate the total gains in market access (MA) from upgrading each link, divided by the cost of upgrading it. The MA of each node is computed as the inverse-travel time weighted sum of GDP of all other nodes, and total MA is the simple sum of MA across all nodes. Total MA is thus denominated in dollars per minute of travel time, and the gains per investment are denominated in dollars per minute per dollar invested. \newline 

Figure \ref{fig:ECAMAPG} shows the return in dollars per minute per dollar invested from building/upgrading each link individually to allow a travel speed of 100$km/h$, while keeping all other links at their original speed. It is thus a partial equilibrium result, as all other links are kept fixed. Ostensibly, re-establishing connectivity along the Russia-Ukraine border yields the highest marginal returns, but also upgrading major transport routes, especially to the major metropoles Moscow and Istanbul, but also to various Western European cities like Rotterdam and Paris, yields high marginal returns. 

\begin{figure}[H]
\centering
\caption{\label{fig:ECAMAPG} Market Access Gains per Dollar Invested}
\includegraphics[width=\textwidth, trim = {0.5cm 0.5cm 0.5cm 0.5cm}, clip]{"../figures/PE/trans_ECA_network_MA_gain_all_100kmh_pusd_yl_or_rd.pdf"}
\includegraphics[width=\textwidth, trim = {0.5cm 0.5cm 0.5cm 0.5cm}, clip]{"../figures/PE/trans_ECA_network_MA_gain_all_100kmh_pusd_cons_MAg0.1.pdf"}
\scriptsize 
\emph{Notes:} Figure shows MA gains in dollars per minute per dollar invested in building/upgrading the respective link. The bottom half highlight high-yield links with return $>$0.1 \$/min/\$ using a stronger colour gradient. 
\end{figure}

The bottom panel of Figure \ref{fig:ECAMAPG} highlights high-yield with a return of $>$0.1 \$/min/\$ using a stronger colour gradient. It reveals that the extremely high-yield segments in red are almost always close to major city centers, at least in Western Europe. In Eastern Europe upgrading the road to Moscow, from Dresden via Warsaw and Minsk, also yields very high returns along extended segments, e.g. between Wroclaw and Lodz in Poland or between Brest and Minsk in Belarus. The other high-return road is the one passing through Eastern Europe down to Istanbul, via Munich, Vienna, Budapest, Belgrade, and Sofia. Here mostly smaller segments, e.g. between Bratislava and Gyor, or in the vicinity of Belgrade, would yield high MA gains if upgraded. \newline 

A caveat of the analysis so far is the unrealistic assumption of open borders throughout the entire region, which is true for the Schengen Area, but not for Eastern Europe or Central Asia. % Thus, additional analysis should take into consideration estimates of border frictions, particularly before prioritizing transnational links. 

\subsubsection*{Optimal Investments in General Equilibrium}

I also consider the optimal general equilibrium spending of a fixed budget on extending and upgrading the network using the quantitative framework of \citet{fajgelbaum2020optimal}. According to my cost estimates, the entire network admits investments of 74.1 billion 2015 USD, of which 51.6 billion fall on upgrading existing links, and 22.5 billion on constructing the 274 proposed links. \newline 

I thus run simulations with investment budgets of 10 and 25 billion to create informative resource allocation problems. In terms of model parameterization, I follow the gridded approach above except that the infrastructure $I_{jk}$ is given in terms of speed ($km/h$) instead of time efficiency, and bounded between the current speed and 100$km/h$. As in \citet{krantz2024optimal}, the estimated per-$km/h$ infrastructure building cost $\delta^I_{jk}$ is the estimated cost of upgrading a link (length times cost/km give in Appendix Figure \ref{fig:ECAcalibdet}) divided by the difference of $I_{jk}$ from 100$km/h$, assuming that investing the estimated amount will bring a link to 100$km/h$. Another difference to the gridded approach are the values of $\gamma$ and $\beta$, which, due to the larger network size, and numerical issues with lower parameter values, are set to $\beta = 1$ and $\gamma = 1.2$ reflecting the more likely increasing returns case for large networks. In particular, setting $\beta = 1$ is crucial to achieve convergence. Despite being a bit ad-hoc, these values are close to \citet{graff2024spatial}, who sets $\gamma = 0.946$ and $\beta = 1.177$. \newline 

The productivity calibration again follows \citet{armington1969theory} and \citet{fajgelbaum2020optimal}: I let 25 key cities produce different products and divide the remaining cities into 4 quartiles by population. Thus, in total there are 29 products traded across the network differentiated by origin. Like \citet{fajgelbaum2020optimal}, I assume that the cross-goods congestion applies. \newline 


Figure \ref{fig:ECAGE10B} shows the optimal \$10B investment allocation. In Western Europe, some funds are spent on building selected new links, whereas in Eastern Europe a combination of new links and upgrades is chosen. Notably, all links connecting Russio to Ukraine are restored, and also a large 


\begin{figure}[H]
\centering
\caption{\label{fig:ECAGE10B} Optimal \$10 Billion Network Investment Allocation}
\includegraphics[width=\textwidth, trim = {0.5cm 0.5cm 0.5cm 0.5cm}, clip]{"../figures/GE/trans_ECA_network_GE_29g_10b_fixed_irs_cgc_sigma3.8_rho0_perc_ug.pdf"}
\includegraphics[width=\textwidth, trim = {0.5cm 0.5cm 0.5cm 0.5cm}, clip]{"../figures/GE/trans_ECA_network_GE_29g_10b_fixed_irs_cgc_sigma3.8_rho0_upw_gain.pdf"}
\scriptsize 
\emph{Notes:} Top panel shows optimal 10 billion 2015 USD network investment in terms of the percentage of proposed construction/upgrading work completed on each link. The bottom panel shows the corresponding welfare gains. \\ \vspace{-5mm}
\end{figure}

\noindent new road connecting St. Petersburg in Russia to Tallinn in Estonia is built. Very little infrastructure on the other hand is built in Central Asia, in fact, only a few links on the east side of the Caspian Sea are upgraded. This is likely due to Central Asia's sparse population and low productivity. \newline 

The bottom panel of Figure \ref{fig:ECAGE10B} provides the welfare gains in each node. The highest gains of up to 2\% are realized in Ukraine in places that are now a war zone. Gains in western Europe are small, generally below 0.5\%, and there are some mild losses, e.g. in central parts of Germany, due to economic activity shifting to the east. It is notably that Central Asia, despite receiving no mentionable investments, still reaps high welfare gains as the investments in Eastern Europe increase connectivity and trade between Western Europe and Central Asia. This is also apparent from the model-implied trade flows, visualized for four key cities in the top panel of Appendix Figure \ref{fig:ECAGETF}: Despite little production and trade from Tashkent, Uzbekistan, the infrastructure investments enable greater trade flows with Moscow, Paris, and Istanbul. 

\begin{figure}[H]
\centering
\caption{\label{fig:ECAGE25B} Optimal \$25 Billion Network Investment Allocation}
\includegraphics[width=\textwidth, trim = {0.5cm 0.5cm 0.5cm 0.5cm}, clip]{"../figures/GE/trans_ECA_network_GE_29g_25b_fixed_irs_cgc_sigma3.8_rho0_perc_ug.pdf"}
\includegraphics[width=\textwidth, trim = {0.5cm 0.5cm 0.5cm 0.5cm}, clip]{"../figures/GE/trans_ECA_network_GE_29g_25b_fixed_irs_cgc_sigma3.8_rho0_upw_gain.pdf"}
\scriptsize 
\emph{Notes:} Top panel shows optimal 25 billion 2015 USD network investment in terms of the percentage of proposed construction/upgrading work completed on each link. The bottom panel shows the corresponding welfare gains. 
\end{figure}

Figure \ref{fig:ECAGE25B} visualizes an optimal \$25B investment. It broadly generates the same spatial patterns of infrastructure allocation and welfare gains, now including full upgrades of many roads in Russia and Eastern Europe and also (partial) upgrades of some Western European roads. \newline 

Appendix Figures \ref{fig:ECAGE10BOU} and \ref{fig:ECAGE25BOU} additionally show optimal \$10B and \$25B investments where the planner is only allowed to upgrade roads. They yield broadly similar allocations as Figures \ref{fig:ECAGE10B} and \ref{fig:ECAGE25B}, with much more intensive upgrading in Eastern Europe and Russia. These should also serve as an important benchmark for policymakers, as most road projects today involve upgrading existing roads rather than fresh construction, and there may be important obstacles to building some of the 274 proposed links in the network other than a rudimentary assessment of geographic feasibility. 


\newpage

\bibliographystyle{apacite}
\bibliography{bibliography} % This links to a file bibliography.bib with the citations

\appendix

\setcounter{figure}{0}
\renewcommand{\thefigure}{A\arabic{figure}}

% \section*{Appendix}

 \begin{figure}[H] % \vspace{-12mm}
 \centering
 \caption{\label{fig:ReallocExtra} Welfare Gains from Optimal Reallocation: Further Countries}
 \vspace{3mm}
 \begin{adjustbox}{center}
 \begin{tabular}{cc}
 Ukraine (0.46\%) & Denmark (0.46\%) \\
 \includegraphics[width=0.6\textwidth, trim = {1cm 3cm 1cm 5cm}, clip]{"../figures/grid_network/realloc_fixed_irs/UKR_opt_realloc_irs.pdf"} &  \includegraphics[width=0.6\textwidth, trim = {1cm 3cm 1cm 5cm}, clip]{"../figures/grid_network/realloc_fixed_irs/DNK_opt_realloc_irs.pdf"} \\ 
 Turkmenistan (0.44\%) & Poland (0.43\%) \\
 \includegraphics[width=0.6\textwidth, trim = {1cm 3cm 1cm 5cm}, clip]{"../figures/grid_network/realloc_fixed_irs/TKM_opt_realloc_irs.pdf"} &  \includegraphics[width=0.6\textwidth, trim = {1cm 3cm 1cm 5cm}, clip]{"../figures/grid_network/realloc_fixed_irs/POL_opt_realloc_irs.pdf"} \\ 
  Greece (0.4\%) & Belarus (0.39\%) \\
 \includegraphics[width=0.6\textwidth, trim = {1cm 3cm 1cm 4cm}, clip]{"../figures/grid_network/realloc_fixed_irs/GRC_opt_realloc_irs.pdf"} &  \includegraphics[width=0.6\textwidth, trim = {1cm 3cm 1cm 4cm}, clip]{"../figures/grid_network/realloc_fixed_irs/BEL_opt_realloc_irs.pdf"} \\ 
 Italy (0.38\%) & Serbia (0.38\%) \\
 \includegraphics[width=0.6\textwidth, trim = {1cm 5cm 1cm 4cm}, clip]{"../figures/grid_network/realloc_fixed_irs/ITA_opt_realloc_irs.pdf"} &  \includegraphics[width=0.6\textwidth, trim = {1cm 3cm 1cm 2cm}, clip]{"../figures/grid_network/realloc_fixed_irs/SRB_opt_realloc_irs.pdf"} 
  \end{tabular}
 \end{adjustbox}
 \end{figure}
 
 
  \begin{figure}[H] \ContinuedFloat \vspace{-5mm}
 \centering
 \caption{\label{fig:ReallocExtra} Welfare Gains from Optimal Reallocation: Further Countries (\emph{Continued})}
 \vspace{3mm}
 \begin{adjustbox}{center}
 \begin{tabular}{cc}
    Turkey (0.38\%) & Azerbaijan (0.38\%) \\
 \includegraphics[width=0.6\textwidth, trim = {3cm 3cm 3cm 7cm}, clip]{"../figures/grid_network/realloc_fixed_irs/TUR_opt_realloc_irs.pdf"} &  \includegraphics[width=0.6\textwidth, trim = {1cm 3cm 1cm 2cm}, clip]{"../figures/grid_network/realloc_fixed_irs/AZE_opt_realloc_irs.pdf"} \\
 France (0.37\%) & Spain (0.36\%) \\
 \includegraphics[width=0.6\textwidth, trim = {1cm 3cm 1cm 2cm}, clip]{"../figures/grid_network/realloc_fixed_irs/FRA_opt_realloc_irs.pdf"} &  \includegraphics[width=0.6\textwidth, trim = {1cm 3cm 1cm 3cm}, clip]{"../figures/grid_network/realloc_fixed_irs/ESP_opt_realloc_irs.pdf"} \\ 
   Uzbekistan (0.34\%) & Romania (0.33\%) \\
 \includegraphics[width=0.6\textwidth, trim = {1cm 3cm 1cm 3cm}, clip]{"../figures/grid_network/realloc_fixed_irs/UZB_opt_realloc_irs.pdf"} &  \includegraphics[width=0.6\textwidth, trim = {1cm 3cm 1cm 3cm}, clip]{"../figures/grid_network/realloc_fixed_irs/ROU_opt_realloc_irs.pdf"} \\ 
 Germany (0.31\%) & Portugal (0.31\%) \\
 \includegraphics[width=0.6\textwidth, trim = {1cm 3cm 1cm 2cm}, clip]{"../figures/grid_network/realloc_fixed_irs/DEU_opt_realloc_irs.pdf"} &  \includegraphics[width=0.6\textwidth, trim = {1cm 3cm 1cm 3cm}, clip]{"../figures/grid_network/realloc_fixed_irs/PRT_opt_realloc_irs.pdf"} \\ 
 \end{tabular}
 \end{adjustbox}
 \end{figure}

 \begin{figure}[H] \vspace{-3mm}
 \centering
 \caption{\label{fig:10PExpExtra} Welfare Gains from Optimal 10\% Expansion: Further Countries}
 \vspace{3mm}
 \begin{adjustbox}{center}
 \begin{tabular}{cc}
 Denmark (0.19\%) & Bulgaria (0.18\%) \\
 \includegraphics[width=0.6\textwidth, trim = {1cm 3cm 1cm 3cm}, clip]{"../figures/grid_network/10perc_fixed_irs/DNK_opt_10perc_irs.pdf"} & \includegraphics[width=0.6\textwidth, trim = {1cm 3cm 1cm 3cm}, clip]{"../figures/grid_network/10perc_fixed_irs/BGR_opt_10perc_irs.pdf"}  \\ 
 Turkmenistan (0.17\%) & Ukraine (0.17\%) \\
 \includegraphics[width=0.6\textwidth, trim = {1cm 3cm 1cm 3cm}, clip]{"../figures/grid_network/10perc_fixed_irs/TKM_opt_10perc_irs.pdf"} & \includegraphics[width=0.6\textwidth, trim = {1cm 3cm 1cm 3cm}, clip]{"../figures/grid_network/10perc_fixed_irs/UKR_opt_10perc_irs.pdf"}  \\ 
  Serbia (0.16\%) & Italy (0.16\%) \\
 \includegraphics[width=0.6\textwidth, trim = {1cm 3cm 1cm 2cm}, clip]{"../figures/grid_network/10perc_fixed_irs/SRB_opt_10perc_irs.pdf"} & \includegraphics[width=0.6\textwidth, trim = {1cm 3cm 1cm 2cm}, clip]{"../figures/grid_network/10perc_fixed_irs/ITA_opt_10perc_irs.pdf"}  \\ 
 Spain (0.15\%) & Turkey (0.14\%) \\
 \includegraphics[width=0.6\textwidth, trim = {1cm 3cm 1cm 3cm}, clip]{"../figures/grid_network/10perc_fixed_irs/ESP_opt_10perc_irs.pdf"} & \includegraphics[width=0.6\textwidth, trim = {3cm 3cm 3cm 3cm}, clip]{"../figures/grid_network/10perc_fixed_irs/TUR_opt_10perc_irs.pdf"}  
 \end{tabular}
 \end{adjustbox}
 \end{figure}
 
 
  \begin{figure}[H]  \ContinuedFloat \vspace{-5mm}
 \centering
 \caption{\label{fig:10PExpExtra} Welfare Gains from Optimal 10\% Expansion: Further Countries (\emph{Continued})}
 \vspace{3mm}
 \begin{adjustbox}{center}
 \begin{tabular}{cc}
    France (0.14\%) & Belarus (0.14\%) \\
 \includegraphics[width=0.6\textwidth, trim = {1cm 2cm 1cm 2cm}, clip]{"../figures/grid_network/10perc_fixed_irs/FRA_opt_10perc_irs.pdf"} & \includegraphics[width=0.6\textwidth, trim = {1cm 3cm 1cm 3cm}, clip]{"../figures/grid_network/10perc_fixed_irs/BLR_opt_10perc_irs.pdf"}  \\ 
 Georgia (0.13\%) & Poland (0.13\%) \\
 \includegraphics[width=0.6\textwidth, trim = {1cm 3cm 1cm 3cm}, clip]{"../figures/grid_network/10perc_fixed_irs/GEO_opt_10perc_irs.pdf"} & \includegraphics[width=0.6\textwidth, trim = {1cm 3cm 1cm 3cm}, clip]{"../figures/grid_network/10perc_fixed_irs/POL_opt_10perc_irs.pdf"}  \\ 
  Albania (0.13\%) & Romania (0.12\%) \\
 \includegraphics[width=0.6\textwidth, trim = {1cm 3cm 1cm 3cm}, clip]{"../figures/grid_network/10perc_fixed_irs/ALB_opt_10perc_irs.pdf"} & \includegraphics[width=0.6\textwidth, trim = {1cm 3cm 1cm 3cm}, clip]{"../figures/grid_network/10perc_fixed_irs/ROU_opt_10perc_irs.pdf"}  \\ 
 Uzbekistan (0.12\%) & Germany (0.12\%) \\
 \includegraphics[width=0.6\textwidth, trim = {1cm 3cm 1cm 3cm}, clip]{"../figures/grid_network/10perc_fixed_irs/UZB_opt_10perc_irs.pdf"} & \includegraphics[width=0.6\textwidth, trim = {1cm 3cm 1cm 2cm}, clip]{"../figures/grid_network/10perc_fixed_irs/DEU_opt_10perc_irs.pdf"}  
 \end{tabular}
 \end{adjustbox}
 \end{figure}
 
 
 \begin{figure}[H]
\centering
\caption{\label{fig:ECAcalibdet} ECA Transport Network Calibration Details}
\begin{tabular}{c}
\includegraphics[width=\textwidth, trim = {0.5cm 0.5cm 0.5cm 0.5cm}, clip]{"../figures/PE/trans_ECA_network_cell_GDPC_const_2017_PPP.pdf"} \\
\includegraphics[width=\textwidth, trim = {0.5cm 0.5cm 0.5cm 0.5cm}, clip]{"../figures/PE/trans_ECA_network_average_link_speed.pdf"} \\
\includegraphics[width=\textwidth, trim = {0.5cm 0.5cm 0.5cm 0.5cm}, clip]{"../figures/PE/trans_ECA_network_type_of_work.pdf"} \\
\includegraphics[width=\textwidth, trim = {0.5cm 0.5cm 0.5cm 0.5cm}, clip]{"../figures/PE/trans_ECA_network_cost_km.pdf"}
\end{tabular}
\end{figure}


\begin{figure}[H]
\centering
\caption{\label{fig:ECAGETF} Model-Generated Trade Flows After Optimal \$10B and \$25B Investments}

After \$10B Optimal Investments \\
\includegraphics[width=\textwidth, trim = {0.5cm 0.5cm 0.5cm 0.5cm}, clip]{"../figures/GE/trans_ECA_network_GE_29g_10b_fixed_irs_cgc_sigma3.8_rho0_good_flows_4_city.pdf"} \\ 
After \$25B Optimal Investments \\
\includegraphics[width=\textwidth, trim = {0.5cm 0.5cm 0.5cm 0.5cm}, clip]{"../figures/GE/trans_ECA_network_GE_29g_25b_fixed_irs_cgc_sigma3.8_rho0_good_flows_4_city.pdf"} 
\end{figure}

\begin{figure}[H]
\centering
\caption{\label{fig:ECAGE10BOU} Optimal \$10 Billion Network Investment Allocation $-$ Only Upgrades}
\includegraphics[width=\textwidth, trim = {0.5cm 0.5cm 0.5cm 0.5cm}, clip]{"../figures/GE/trans_ECA_network_GE_ug_29g_10b_fixed_irs_cgc_sigma3.8_rho0_perc_ug.pdf"}
\includegraphics[width=\textwidth, trim = {0.5cm 0.5cm 0.5cm 0.5cm}, clip]{"../figures/GE/trans_ECA_network_GE_ug_29g_10b_fixed_irs_cgc_sigma3.8_rho0_upw_gain.pdf"}
\includegraphics[width=\textwidth, trim = {0.5cm 0.5cm 0.5cm 0.5cm}, clip]{"../figures/GE/trans_ECA_network_GE_ug_29g_10b_fixed_irs_cgc_sigma3.8_rho0_good_flows_4_city.pdf"}
\scriptsize 
\emph{Notes:} Top panel shows optimal 10 billion 2015 USD network investment in terms of the percentage of proposed upgrading completed. Middle panel shows the corresponding welfare gains. Bottom panel the generated trade flows. 
\end{figure}


\begin{figure}[H]
\centering
\caption{\label{fig:ECAGE25BOU} Optimal \$25 Billion Network Investment Allocationn $-$ Only Upgrades}
\includegraphics[width=\textwidth, trim = {0.5cm 0.5cm 0.5cm 0.5cm}, clip]{"../figures/GE/trans_ECA_network_GE_ug_29g_25b_fixed_irs_cgc_sigma3.8_rho0_perc_ug.pdf"}
\includegraphics[width=\textwidth, trim = {0.5cm 0.5cm 0.5cm 0.5cm}, clip]{"../figures/GE/trans_ECA_network_GE_ug_29g_25b_fixed_irs_cgc_sigma3.8_rho0_upw_gain.pdf"}
\includegraphics[width=\textwidth, trim = {0.5cm 0.5cm 0.5cm 0.5cm}, clip]{"../figures/GE/trans_ECA_network_GE_ug_29g_25b_fixed_irs_cgc_sigma3.8_rho0_good_flows_4_city.pdf"}
\scriptsize 
\emph{Notes:} Top panel shows optimal 25 billion 2015 USD network investment in terms of the percentage of proposed upgrading completed. Middle panel shows the corresponding welfare gains. Bottom panel the generated trade flows. 
\end{figure}
 
 
\end{document}